\documentclass[10pt,a4paper]{amsart}
\usepackage[margin=19mm]{geometry}
\usepackage{amsmath,amssymb,mathtools}
\usepackage[scr=rsfs,cal=euler]{mathalfa}
\usepackage{xfrac}
\usepackage[unicode,hidelinks,bookmarks=false]{hyperref}
\newcommand{\ie}{i.e.\ }
\newcommand{\cf}{cf.\ }
\newcommand{\resp}{respectively\ }
\newcommand{\Subsection}{\subsection}
\providecommand{\nobreakdash}{\nobreak\hbox{-}}
\setlength{\emergencystretch}{2em}
\multlinegap0pt
\usepackage{bm}
\usepackage{bbm}
\usepackage{dsfont}
\usepackage{xspace}
\usepackage{cancel}
\usepackage{mathtools}
\usepackage{amsthm}
\makeatletter
\@ifundefined{defn}{\newtheorem{defn}{Defn}}{}
\@ifundefined{thm}{\newtheorem{thm}[defn]{Theorem}}{}
\makeatother
\title[Upper bounds for the list-distinguishing chromatic number]{Upper bounds for the list-distinguishing chromatic number}
\author{Amitayu Banerjee, Zal\'an Moln\'ar, Alexa Gopaulsingh}
\date{Source: arXiv:2405.12733v3 (CC BY 4.0)}
\begin{document}
\maketitle

\noindent\emph{Source and licence.} Upper bounds for the list-distinguishing chromatic number. Version arXiv:2405.12733v3 (CC BY 4.0), distributed under the Creative Commons Attribution 4.0 International licence, as recorded in the arXiv licence metadata for this exact version and shown on the abstract page https://arxiv.org/abs/2405.12733v3. Full rights notice: licenses/arxiv-2405.12733-CC-BY-4.0.txt.\par\smallskip

\noindent\textit{Editorial scope.} Counted unit: the Thm whose source label is \texttt{Theorem 3.1} in the article (source0.tex lines 419--421, source label \texttt{Theorem 3.1}), delivered together with the article's own complete original proof of it (source0.tex lines 423--488, one \texttt{proof} environment of 66 lines). No mathematics is written, repaired or re-derived: statement and proof are the article's own text line for line. Required context retained uncounted: Theorem 2.7 (lines 385--407). Every definition, display and label of those lines is retained unchanged. Omitted from this package and not redistributed: 1--384, 408--418, 422--422, 489--1165. Nothing inside a retained excerpt is omitted or altered: every retained body is delivered exactly as the article's lines are written, so no omission receipt of any kind accompanies this package. Citations: the retained excerpts carry no citation command at all, so no bibliography entry of the article is transcribed and no reference is redistributed. Reference adaptation: none was needed and none was made; every reference inside the delivered unit resolves inside it, so no unresolved reference or citation is produced. Printed slips kept literal: no printed word, symbol, hypothesis or number of the counted statement or of its proof was altered. Layout and class: the article's own class is \texttt{amsart} with packages including hyperref, enumerate, geometry, amsmath, amssymb, amsthm, amsfonts, graphicx, tikz; the wrapper is the lane's pinned \texttt{amsart} editorial wrapper at 10pt on A4, which restates the theorem-like environments the retained text needs and adds the article's own macro definitions that the retained text needs (none), with extra wrapper packages (bm, bbm, dsfont, xspace, cancel, mathtools). The delivered numbering is the unit's own. Assets: no retained span contains a figure environment, a graphics-inclusion command, a PDF-page inclusion, a table or a TikZ picture; no image file is copied into this package, the package declares no asset dependency and no image file is redistributed with it. Licence: version arXiv:2405.12733v3 of this article is distributed under the Creative Commons Attribution 4.0 International licence, as recorded in the arXiv licence metadata for this exact version and shown on the abstract page https://arxiv.org/abs/2405.12733v3. A full rights notice is delivered at licenses/arxiv-2405.12733-CC-BY-4.0.txt. Characters: every retained excerpt is pure ASCII, so the delivered engine, XeLaTeX with its default Latin Modern text font, reports no missing character.
\begin{thm}\label{Theorem 2.7}
    {\em The entries in the following table are correct.}
    
\begin{center}
      \begin{tabular}{|ll|l|l|}
        \hline
             & Graph & $\chi_{D_L}(G)$ \\
             \hline\hline
             & $P_{2t}, t\geq 1$ & $2$\\
             & $P_{2t+1}, t\geq 1$ & $3$ \\
             & $C_4$ & $4$ \\
             & $C_5$ & $3$ \\
             & $C_6$ & $4$\\
             & $C_{2m+1}, m\geq 3$  & $3$ \\
             & $C_{2n}, n\geq 4$ & $3$ \\
               
        \hline
        \end{tabular}
        \\
       Table 1. 
        \label{tab:my_label}
    \end{center}  
\end{thm}

\begin{thm}\label{Theorem 3.1}
{\em Let $G$ be a connected finite graph with maximum degree $\Delta(G)$. Then $\chi_{D_{L}}(G)\leq 2\Delta(G)-1$, unless $G$ is  $K_{\Delta(G),\Delta(G)}$ or $C_6$. In these cases, $\chi_{D_{L}}(G) = 2\Delta(G)$}.    
\end{thm}

\begin{proof}
   If $\Delta(G)=2$, then by Theorem \ref{Theorem 2.7}, we have $\chi_{D_L}(G) = 4$ if and only if $G$ is $C_4$ or $C_6$. It is easy to see that $\chi_{D_L}(K_{\Delta(G),\Delta(G)}) = 2\Delta(G)$, and $K_{2,2}$ is $C_4$. Hence, we may assume that $\Delta(G) \geq 3$ and $G$ is not $K_{\Delta(G),\Delta(G)}$. We show that in this case $\chi_{D_L}(G) \leq 2\Delta(G)-1$. Let $L=\{L(v)\}_{v\in V_G}$ be an assignment of lists with $|L(z)| = 2\Delta(G)-1$, for all $z\in V_G$. 
   Let $v\in V_G$ be a vertex with degree $\Delta(G)$ and $T$ be a breadth-first search ($BFS$) spanning tree of $G$ rooted at $v$. We use the notation $<$ to denote the $BFS$ order. Also, for $z\in V_{G}$, by $N(z)$ we denote the set of its neighbors in $G$ and $S(z)$ denotes the set of its siblings in $T$. We will define partial functions $f_x: V_{G} \to \bigcup_{z\in V_{G}}L(z)$ for $x\in V_{G}$ such that the following holds: 
\begin{enumerate}
    \item[(i)] if $x < y$, then $f_x\subseteq f_y$,
    \item[(ii)] if $f=\bigcup_{x\in V_{G}} f_x$, then $f(v)\in L(v)$ for all $v\in V_{G}$,
    
    \item[(iii)] the domain of $f_{x}$ includes all vertices till $x$ in the $BFS$ order.
\end{enumerate}
For the base case, we assign pairwise different colors to $v$ and to members of $N(v) = \{v_1,\dots, v_{\Delta(G)}\}$ from the respective lists, say $c_v\in L(v)$ and $c_{v_i}\in L(v_i)$, where $1\leq i\leq \Delta(G)$. 
For the rest, we try to keep $v$ to be the only vertex of $G$ with the property of being colored with $c_v$, while all its neighbors are colored by $c_{v_1},\dots, c_{v_{\Delta(G)}}$. We will refer to this property as ($\ast$). 

We proceed by coloring the vertices of $G$ in their $BFS$ order inductively, as follows: Let $x\in V_{G}$ be the $<$-minimal element for which no color has been assigned and $z<x$ be the immediate $BFS$ predecessor of $x$. We define $f_x$. Let us write,
$$
    S_<(x) = \{y\in S(x): y<x\} \text{ and } N_<(x) = \{y\in N(x): y<x\}.
$$
Then $f_z$ has already assigned colors to each member of $S_<(x) \cup N_<(x)$. Define,
$$
A_x= L(x)\setminus \Big(\{c_v\} \cup f_z[S_<(x)] \cup f_z[N_<(x)]\Big).
$$
If $A_{x} \neq \emptyset$, then let $f_x = f_z \cup \{\langle x, c \rangle\} $,
where $c\in A_x$. If $A_x=\emptyset$, then since $|N(x)| \leq \Delta(G)$ and $|S(x)|\leq \Delta(G)-2$, we must have used $|S_<(x)| + |N_<(x)| = 2\Delta(G) -2$ many colors from $L(x)$ to color the neighbors and siblings that come before $x$. Thus, the following holds: 

\begin{enumerate}
    \item[$(a)$] $c_v\in L(x)$,
    \item[$(b)$] $S_<(x)=S(x)$ and $N_<(x)=N(x)$,
    \item[$(c)$] $|S(x)| = \Delta(G) -2$ and $|N(x)| = \Delta(G)$,
    \item[$(d)$] $N(x) \cap S(x) = \emptyset$,
    \item[$(e)$] for all $y\in S(x) \cup N(x)$, we have $f_z(y)\in L(x)$.
\end{enumerate}
In this case, let $f_x = f_z\cup \{\langle x, c_v \rangle\}$.
In either case, $f_x$ has (i) and (iii). 
Finally, set $f=\bigcup_{x\in V_{G}}f_x$. It is clear that $f$ satisfies (ii), and whenever $f(x)=c_v$ where $x\neq v$, then $(a)-(e)$ must hold. Moreover, if no vertex other than $v$ has property $(\ast)$, then it is easy to see that $f$ is a proper distinguishing coloring. 
However, suppose there is a vertex other than $v$, say $x$, with property $(\ast)$ and is the $<$-minimal such vertex. Then $(a)-(e)$ are satisfied for $x$. We modify $f$ in such a way that $x$ has no longer $(\ast)$, while it remains a proper list coloring.
By a careful introspection we will analyze the following cases: \footnote{We note that \textsc{Case 1} and {\em Case A} of \textsc{Subcase 3.1} arise when the lists are nonidentical.}

\noindent \textsc{Case 1:} If there is $w\in S(x)$ such that $|N(w)| \neq |N(x)|$, then color $x$ with $f(w)$,
which can be done by $(e)$. Then there is no color-preserving automorphism of $G$, which fixes $v$ and maps $x$ to $w$. Meanwhile, by $(d)$, this keeps the coloring proper, and $x$ has no longer the property $(\ast)$.


\noindent \textsc{Case 2:} For every $z\in S(x)$, we have $|N(z)| = |N(x)|$, and there is a $w\in S(x)$ such that $N(w)\neq N(x)$. By $(c)$, fix $y\in N(x)$ such that $y\not\in N(w)$. Since the neighbors of $x$ must have occurred before $x$ in the $BFS$ ordering, we have $f(w), f(y)\in L(x)$ by $(e)$.
We color $x$ with $f(w)$. The rest follows the arguments of \textsc{Case 1}.

\noindent  \textsc{Case 3: } For every $w\in S(x)$, we have $N(w) = N(x)$.
\begin{enumerate}
    \item[]\textsc{Subcase 3.1:} There is a vertex $z\in N(x)$ with no sibling or parent colored with $c_v$. We need to consider two different cases.
    \begin{enumerate}
    \item[]\textit{Case A:}
    If $c_v\not\in L(z)$, then by $(c)$ at most $(\Delta(G)-1) + (\Delta(G)-2) = 2\Delta(G)-3$ number of colors can appear in the list $L(z)$ that are used for coloring $N(z)\cup S(z)$. Hence, there must be $c\in L(z)$ such that $c\neq f(z)$, which also differs from all the colors assigned to $N(z)\cup S(z)$. 
    We color $z$ with $c$ and $x$ with $f(z)$.
    %Redefine $f$ by  $z\mapsto c$.
    Then $x$ no longer has $(\ast)$ (as no member of $N(v)$ has color $c$), and this modification keeps the coloring proper.
    Moreover, $z$ does not have the property 
    $(\ast)$ as $z$ is colored by $c$, but $c\neq c_{v}$, since $c\in L(z)$ and $c_{v}\not\in L(z)$.
   \item[]  \textit{Case B:} If $c_v\in L(z)$, then color $x$ with $f(z)$ and $z$ with $c_{v}$, which can be done by $(e)$. Clearly, $x$ does not have $(\ast)$, and $z$ does not have $(\ast)$ either by $(b)-(d)$. 
    \end{enumerate}
   
    \item[] \textsc{Subcase 3.2:} Each $y \in N(x)$ has a sibling or parent colored $c_v$. By $(b)$, each $y\in N(x)$ and its siblings must have come before $x$ in the $BFS$ order. Hence, their parent cannot have ($\ast$), unless it is the root $v$. There must exist a $z\in N(x)$ such that $z$ is not a child of $v$. Otherwise, $G$ is $K_{\Delta(G),\Delta(G)}$, which contradicts the assumption that $G$ is not $K_{\Delta(G),\Delta(G)}$. 
    Fix such $z$.
    Let $B(z)$ denotes the set of colors assigned to the vertices of $N(z)\cup S(z)$. We claim that the set $L'(z)=L(z)\backslash \{f(z)\cup B(z)\}$ is non-empty. Let $p$ be the parent of $z$ in $T$. Then $N(z) = S(x) \cup \{p\}\cup \{x\}$, $\vert S(z)\vert \leq \Delta(G)-2$, and $\vert N(z)\vert = \vert S(x)\cup \{p\}\cup \{x\}\vert=\Delta(G)$. The color $c_v$ was used once for $x$ and once among the vertices of $S(z) \cup \{p\}$. Thus, $\vert B(z)\vert\leq 2\Delta(G)-3$. Since $\vert L(z)\vert= 2\Delta(G)-1$, the set $L'(z)$ is non-empty.
    Pick $c\in L'(z)$.
    Using $(e)$, we can color $x$ with $f(z)$ and $z$ with $c$.  We note that $z$ does not get property $(\ast)$. This is because $c\neq c_{v}$ since $c_{v} \in B(z)$.
    Similarly, as in the previous cases, the new coloring remains a proper list coloring, and $v$ is the only $<$-predecessor of $x$ having $(\ast)$, and all of them are fixed as long as $v$ is fixed.
\end{enumerate}
By iterating this procedure, we can recolor every node with property $(\ast)$, and hence we obtain a list-distinguishing proper coloring of $G$.
\end{proof}

\end{document}
